\section{Твердое тело}
\subsection{Равновесие палочки}
%%%%%%%%%%%%%%%%%%%%%%%%%%%%%%%%%%%%%%%%%%%%%%%%%%%%
%%%%%%%%%%%%%%%%%%%%%%%%%%%%%%%%%%%%%%%%%%%%%%%%%%%%
 \z[{\begin{wrapfigure}[6]{r}{0.3\textwidth}%
		\vspace{-12pt}%
		\begin{flushright}%
 			\includegraphics[width=0.3\textwidth]{8.1.1}%
		\end{flushright}%
	\end{wrapfigure}%
	}]%
	{На двух вертикально расположенных пружинах, в недеформированном состоянии имеющих одинаковую длину, горизонтально подвешен невесомый стержень. Жесткости пружин $k_1 = 0.02$ Н/м и $k_2 = 0.03$ Н/м, расстояние между ними $d = 1$ м. На каком расстоянии $x$ от первой пружины нужно подвесить груз, чтобы стержень остался в горизонтальном положении?}
%
{ $x = \frac{d \cdot k_2}{k_1 + k_2} = 0.6$ м.}
%%%%%%%%%%%%%%%%%%%%%%%%%%%%%%%%%%%%%%%%%%%%%%%%%%%%
\z[{\begin{wrapfigure}[5]{r}{0.3\textwidth}%
		\vspace{-15pt}%
		\begin{flushright}%
 			\includegraphics[width=0.3\textwidth]{8.1.2}%
		\end{flushright}%
	\end{wrapfigure}%
	}]%
	{Линейка массы $0.01$ кг лежит на двух опорах так, как это показано на рисунке. На один конец линейки положен груз. Какова масса груза, при которой возможно равновесие?}
%
{$m = 7.5$ г.}
%%%%%%%%%%%%%%%%%%%%%%%%%%%%%%%%%%%%%%%%%%%%%%%%%%%% 
 \z[{\begin{wrapfigure}[3]{r}{0.3\textwidth}%
		\vspace{-15pt}%
		\begin{flushright}%
 			\includegraphics[width=0.3\textwidth]{8.1.3}%
		\end{flushright}%
	\end{wrapfigure}%
	}]%
	{Стержень массой $m$ и длиной $l$ удерживают в горизонтальном положении с помощью двух точечных опор, расположенных на расстоянии $l/5$ друг от друга. Найти силы реакции опор, считая, что на одну из них стержень опирается самым краем.}
%
{$N_1 = \frac{5}{2}mg$, $N_2 = \frac{3}{2}mg$}
%%%%%%%%%%%%%%%%%%%%%%%%%%%%%%%%%%%%%%%%%%%%%%%%%%%% 
 \z{Под каким минимальным углом $\alpha$ к горизонту может стоять лестница, прислоненная к вертикальной стене, если ее центр масс находится в середине? Коэффициент трения между лестницей и полом равен $\mu_1$, а между лестницей и стеной $\mu_2$.}
%
{$\varphi = \arctg\left(\frac{1 - \mu_1 \mu_2}{2\mu_2}\right)$}
%%%%%%%%%%%%%%%%%%%%%%%%%%%%%%%%%%%%%%%%%%%%%%%%%%%%
%%%%%%%%%%%%%%%%%%%%%%%%%%%%%%%%%%%%%%%%%%%%%%%%%%%%	
 \z[{\begin{wrapfigure}[4]{r}{0.3\textwidth}%
		\vspace{-15pt}%
		\begin{flushright}%
 			\includegraphics[width=0.3\textwidth]{8.1.5}%
		\end{flushright}%
	\end{wrapfigure}%
	}]%
	{В гладкой закрепленной полусфере свободно лежит палочка массы $m$ так, что угол ее с горизонтом равен $\alpha$, а конец выходит за край полусферы. С какими силами действует палочка на полусферу в точках соприкосновения $A$ и $B$?}
 %
{$N_A = Mg \tan \alpha$, $N_B = Mg \frac{\cos 2\alpha}{\cos \alpha}$}